\section{A pressure law for a signed drift}\label{sec:pressure}

The local metric enlargement will be obtained from a graph equation whose
right-hand side depends on a height discrepancy and on the drift.  We need
this dependence to be increasing in the discrepancy, with poles that confine
the graph height.  We also need quantitative control of its drift derivative
for every real value of the drift.  The following construction adapts the
pressure method of \cite[Theorem~2.1]{OpenAIWidth} to this signed parameter.

\begin{proposition}[Signed pressure law]\label{pressure:law}
For every \(A_0,P_0,\epsilon,b_0>0\), there are constants
\(S_0\in(0,1]\), \(T\ge10\), and a smooth function
\[
 F:(-S_0,S_0)\times\R\longrightarrow\R
\]
that is odd in its first variable and has the following properties.  Set
\begin{equation}\label{pressure:reserve}
 N(t)=\sqrt{t^2+T^2},\qquad V(t)=N(t)^2=t^2+T^2.
\end{equation}
At every \((x,t)\in(-S_0,S_0)\times\R\),
\begin{equation}\label{pressure:bounds}
 \begin{gathered}
 F_x\ge A_0N(t),\qquad |F_t|\le\epsilon|F|,
       \qquad FF_t\le\epsilon F_x,\\
 F_x\le b_0N(t)^2\quad\text{or}\quad -FF_t\ge P_0F_x,\\
 F_x\le b_0N(t)^2+\epsilon F^2.
 \end{gathered}
\end{equation}
Moreover, \(F(x,t)\to+\infty\) as \(x\uparrow S_0\) and
\(F(x,t)\to-\infty\) as \(x\downarrow-S_0\), uniformly for \(t\)
in each compact subset of \(\R\).
\end{proposition}

In the local scalar-curvature comparison, the alternative permits a
negative term proportional to \(F_x\) to be controlled either by a small
multiple of the reserve \(V(t)\) or by a positive contribution from
\(-FF_t\).  Where the graph is nearly horizontal, the comparison also
contains a positive multiple of \(F^2\).  The last inequality lets us
absorb the \(\epsilon F^2\) part of the derivative loss into this positive
term and reduce the remaining reserve loss by making the metric
modification small there.

\begin{proof}
It is easier to prescribe the discrepancy \(x=X(f,t)\) as a function
of a positive pressure value \(f\), and then invert.  Indeed, the implicit
relation \(X(F(x,t),t)=x\) gives
\[
 F_x=\frac1{X_f},\qquad
 \frac{FF_t}{F_x}=-fX_t.
\]
Thus small positive \(X_f\) gives a large height derivative, while the sign
of \(X_t\) controls the mixed term.  We will arrange two regions: one where
\(F_x\le b_0N^2\), and one where \(-FF_t\ge P_0F_x\).

\paragraph{Choosing the inverse form.}
Choose
\[
 k\ge\max\{A_0,\epsilon^{-1},\pi/2\},\qquad
 Y\ge1,\qquad Y^{-3}\le\frac{b_0}{12k}.
\]
For now let \(T\ge10\) be arbitrary.  The simple candidate
\(F(x,t)=N(t)\tan(kx)\) has the desired poles.
However, for fixed \(t>0\), its derivative \(F_x/N(t)^2\) tends to
infinity near a pole while \(FF_t>0\).  It therefore satisfies neither
part of the required alternative there.  We will retain this candidate
at small pressure values and change it at large ones to make \(F_t<0\)
when \(F>0\).

Given \(f>0\) and \(t\in\R\), set
\[
 q=f^{2/3},\qquad y=\frac tq,\qquad z=\frac Tq,
       \qquad R=\sqrt{y^2+z^2}=\frac{N(t)}q.
\]
For fixed \(t\), increasing \(f\) moves \((y,z)\) toward the origin
along a ray; for fixed \(f\), varying \(t\) moves it horizontally.
For a positive smooth function \(G\) on \(\R^2\), to be chosen below,
consider the inverse expression
\begin{equation}\label{pressure:inverse}
 X(f,t)=\frac1k\arctan\!\left(\frac{\sqrt q}{G(y,z)}\right),
       \qquad S_0=\frac\pi{2k}\le1.
\end{equation}
The identity \(q^{3/2}=f\) ensures that, wherever \(G=R\), we have
\(\sqrt q/G=f/N(t)\), recovering the tangent candidate.
Differentiating the inverse expression identifies the properties we need
from \(G\):
\begin{equation}\label{pressure:inverse-derivatives}
 \begin{gathered}
 A_G=G+2(yG_y+zG_z),\\
 fX_t=-\frac{G_y}{k(1+G^2/q)},\qquad
 X_f=\frac{A_G}{3kq^2(1+G^2/q)}.
 \end{gathered}
\end{equation}
These identities follow by differentiating \(q=f^{2/3}\), using
\(y_q=-y/q\), \(z_q=-z/q\) for the second one.  Thus \(A_G>0\)
will make \(X\) increasing in \(f\), while \(G_y<0\) will give the
favorable mixed sign.  Taking \(G=R\) outside a compact set recovers the
tangent candidate at small \(f\).  The modification near the origin
will control large \(f\).

\paragraph{Constructing the auxiliary function.}
We now construct \(G\) on the whole \((y,z)\)-plane, with
\(R=\sqrt{y^2+z^2}\).  It will equal \(R\) outside a compact set and
be strictly decreasing in \(y\) on the disk \(R\le Y\).  Define
\[
 m(y)=D_G\exp\bigl(-\arctan(y)/8\bigr),
\]
where \(D_G>0\) will be fixed shortly.  Choose a nonnegative smooth
probability density \(\varrho\) with compact support in \((1,2)\), and set
\begin{equation}\label{pressure:rounded-max}
 G(y,z)=\int_1^2\max\{R,\lambda m(y)\}\varrho(\lambda)\,\dd\lambda.
\end{equation}
To see that this averaging produces a smooth function, let
\[
 \Phi(v)=\int_1^2\max\{v,\lambda\}\varrho(\lambda)\,\dd\lambda.
\]
Then \(\Phi''=\varrho\), with the density extended smoothly by zero,
and \(G=m\Phi(R/m)\).  This proves smoothness away from the origin.
Near the origin every integrand in \eqref{pressure:rounded-max} equals
\(\lambda m\), so \(G\) is smooth there as well.  Since \(m\) is bounded
above and below by positive constants, \(G\ge R\), \(G>0\), and
\(G=R\) outside a compact set.

The derivatives needed for inversion are
\begin{equation}\label{pressure:G-bounds}
 G_y\le1,\qquad \frac G2\le A_G\le3G.
\end{equation}
Here \(m'=-m/(8(1+y^2))\), so \(|ym'|\le m/16\).
On the radial branch \(R\), the corresponding expressions are
\(R_y\le1\) and \(R+2(yR_y+zR_z)=3R\).
On the other branch \(\lambda m\), the first derivative is nonpositive
and
\[
 \frac78\lambda m
 \le\lambda(m+2ym')\le\frac98\lambda m.
\]
Averaging these branchwise inequalities gives \eqref{pressure:G-bounds};
the equality set of the two branches has measure zero in \(\lambda\)
at each point away from the origin.

Let \(c_{\varrho}=\int_1^2\lambda\varrho(\lambda)\,\dd\lambda\), so
\(1<c_{\varrho}<2\).  Increase \(D_G\) until, for \(|y|\le Y\),
\[
 m(y)>Y,\qquad m(y)\ge16kP_0(1+Y^2).
\]
On the disk \(R\le Y\), we then have \(G=c_{\varrho}m\) and
\begin{equation}\label{pressure:G-inner}
 G_y=-\frac{c_{\varrho}m}{8(1+y^2)}\le-2kP_0.
\end{equation}
The function \(G\), as well as \(k\) and \(Y\), depends only on the four
given constants and not on \(T\).

\paragraph{Inversion and the basic bounds.}
For every \(T\ge10\), the function \(G\) just constructed gives
\(X_f>0\) by \eqref{pressure:inverse-derivatives} and
\eqref{pressure:G-bounds}.
For fixed \(t\), as \(f\downarrow0\) we eventually have \(G=R=N(t)/q\),
and hence
\begin{equation}\label{pressure:near-zero}
 X(f,t)=\frac1k\arctan\!\left(\frac f{N(t)}\right).
\end{equation}
As \(f\to\infty\), the point \((y,z)\) tends to the origin and
\(\sqrt q/G(y,z)\to\infty\).  Therefore \(X(\cdot,t)\) increases
smoothly from \(0\) to \(S_0\).  The implicit-function theorem and
\(X_f>0\) give a jointly smooth inverse \(F(x,t)>0\) for
\(0<x<S_0\).  Formula
\eqref{pressure:near-zero} holds uniformly for \(t\) in compact ranges;
its inverse is \(F(x,t)=N(t)\tan(kx)\) near \(x=0\).
We may thus extend \(F\) smoothly by oddness across \(x=0\).

For completeness, the pole assertion is also uniform in the parameter.
Given a compact set \(K\subset\R\) and \(f_0>0\), continuity gives
\(\max_{t\in K}X(f_0,t)<S_0\).  For \(x\) above this maximum,
monotonicity implies \(F(x,t)>f_0\) for every \(t\in K\).
Oddness gives the assertion at the other endpoint.

It remains to choose \(T\) so that this smooth inverse satisfies all the
quantitative bounds.  At \(x>0\), with \(f=F(x,t)\),
\begin{equation}\label{pressure:implicit-derivatives}
 \frac{FF_t}{F_x}=-fX_t,\qquad
 \frac{F_t}{F}=\frac{3G_y}{qA_G}.
\end{equation}
By \eqref{pressure:G-bounds} and \(R\le G\),
\[
 N(t)X_f=\frac{RA_G}{3k(q+G^2)}\le\frac1k,
 \qquad -fX_t\le\frac1k.
\]
Thus \(F_x\ge A_0N(t)\) and \(FF_t\le\epsilon F_x\).
The finite constant
\[
 C_G=\sup_{(y,z)\in\R^2}\frac{R|G_y|}{A_G}
\]
controls the remaining drift derivative: it is finite on compact sets by
positivity of \(A_G\), and outside a compact set we have \(G=R\) and
\(A_G=3R\).  Consequently,
\begin{equation}\label{pressure:t-derivative}
 \left|\frac{F_t}{F}\right|
 \le\frac{3C_G}{N(t)}\le\frac{3C_G}{T}.
\end{equation}

\paragraph{The two regions and the final choice of \(T\).}
For \(R>Y\), let \(B=\sup_{R>Y}G/R<\infty\).  Since \(q=N(t)/R\ge T/R\),
\eqref{pressure:inverse-derivatives} gives
\[
 N(t)^2X_f
 =\frac{R^2A_G}{3k(1+G^2/q)}
 \ge\frac{R^3}{6k(1+B^2R^3/T)}.
\]
In particular,
\begin{equation}\label{pressure:outer-bound}
 \frac{F_x}{N(t)^2}
 \le6k\left(R^{-3}+\frac{B^2}{T}\right)
 \le6k\left(Y^{-3}+\frac{B^2}{T}\right).
\end{equation}
This will be at most \(b_0\).

The disk \(R\le Y\) corresponds to
\(f\ge(N(t)/Y)^{3/2}\), the large-pressure region.  Set
\[
 M_G=\max_{R\le Y}G,\qquad a_G=\min_{R\le Y}A_G>0.
\]
Here \(q\ge T/Y\), so choosing \(T\ge YM_G^2\) ensures
\(G^2/q\le1\).  Combining this with \eqref{pressure:G-inner} yields
\[
 fX_t\ge P_0,
 \qquad
 f^2X_f=\frac{qA_G}{3k(1+G^2/q)}
       \ge\frac{Ta_G}{6kY}.
\]
The first inequality is exactly \(-FF_t\ge P_0F_x\); the second will
give \(F_x\le\epsilon F^2\).

All constants in these conditions were fixed before \(T\).  We now choose
\begin{equation}\label{pressure:T-choice}
 T\ge\max\left\{10,\frac{3C_G}{\epsilon},
       \frac{12kB^2}{b_0},\ YM_G^2,
       \frac{6kY}{\epsilon a_G}\right\}.
\end{equation}
Then \eqref{pressure:t-derivative} proves \(|F_t|\le\epsilon|F|\),
and \eqref{pressure:outer-bound}, together with the choice of \(Y\),
proves \(F_x\le b_0N(t)^2\) when \(R>Y\).
On \(R\le Y\), we have both \(-FF_t\ge P_0F_x\) and
\(F_x\le\epsilon F^2\).  These are precisely the asserted alternative,
and in either region they also imply the last inequality in
\eqref{pressure:bounds}.  Finally, under \(x\mapsto-x\), both \(F\) and
\(F_t\) change sign while \(F_x\) does not.  All the bounds therefore
hold for negative \(x\), and continuity gives them at zero.  No step
restricted the sign of \(t\).
\end{proof}
